% Gerado por avaliacao/relatorio.py — nao editar a mao; rode novamente a avaliacao.
\begin{table}[htbp]
\caption{Portão de suficiência: separação entre demandas e pedidos fora do domínio}\label{tab:portao}
\centering\scriptsize
\setlength{\tabcolsep}{3.5pt}
\resizebox{\linewidth}{!}{%
\begin{tabular}{lrrlrrrl}
\toprule
Pontuação & AUC & $\rho$ & v0.2.0: TPR / FPR & Limiares (VC) & TPR / FPR (VC) & IC 95\% FPR (VC) & Adotado: TPR / FPR \\
\midrule
hash & 0,983 & 0,38 & 0,10: 100\% / 100\% & 0,49 e 0,49 & 97\% / 6\% & 2--16 & 0,46: 97\% / 10\% \\
spaCy & 0,923 & 0,13 & -- & 0,77 e 0,73 & 84\% / 26\% & 16--40 & 0,74: 91\% / 26\% \\
BGE-M3 & 0,934 & 0,05 & 0,45: 100\% / 62\% & 0,53 e 0,55 & 81\% / 12\% & 6--24 & 0,52: 94\% / 22\% \\
BM25 & 0,983 & 0,03 & -- & 0,52 e 0,52 & 97\% / 6\% & 2--16 & não se aplica \\
\bottomrule
\end{tabular}}
\fonte{Elaborado pelo autor, a partir de \texttt{avaliacao/resultados/resumo.json}.}
\nota{32 demandas e 50 pedidos fora do domínio. TPR: demandas aceitas; FPR: pedidos fora do domínio aceitos. VC: validação cruzada em duas partes (limiar escolhido numa metade pelo índice de Youden e aplicado à outra). $\rho$: correlação de Spearman entre o tamanho da consulta e a pontuação, entre as demandas. IC 95\%: intervalo de Wilson, em pontos percentuais.}
\end{table}
